\uuid{HRKt}
\titre{Pgcd}
\auteur{bodin ; croizat ; sacre}
\organisation{exo7}
\licence{CC BY-NC-SA 4.0}
\datecreate{2022-03-08}
\niveau{L1}
\module{Algèbre}
\chapitre{Arithmétique dans Z}
\sousChapitre{Pgcd, ppcm, algorithme d'Euclide}
\difficulte{2}
\theme{pgcd}
\exo7qcmid{hilisit.arithmetique.009}

\contenu{
\begin{question}
\qclassification{103.03}
    Pour $a=24$ et $b=8$, on a :
    \begin{answers}
        \bad{$\ppcm(a,b) = 8$.}
        \good{$\ppcm(a,b) = 24$.}
        \good{$a$ est un multiple de $b$.}
        \bad{$a$ est dans la liste des diviseurs de $b$.}
    \end{answers}
    \begin{explanations}
     $a$ étant un multiple de $b$ (on a $24 = 8 \times 3$), on a immédiatement $\pgcd(a,b) = b$ et $\ppcm(a,b) = a$.
    \end{explanations}
\end{question}
}
